\chapter{Ưu Điểm Có Thật, Nhưng Không Cứu Nổi Cách Dùng Tệ}
\markboth{Ưu Điểm Có Thật, Nhưng Không Cứu Nổi Cách Dùng Tệ}{Ưu Điểm Có Thật, Nhưng Không Cứu Nổi Cách Dùng Tệ}

\begin{center}
	\textit{\color{bookprimary}``Lợi ích của điện thoại là thật. Nhưng rất nhiều người đang lấy vài hạt đường để bào chữa cho cả một đống rác họ nuốt vào mỗi ngày.''}

	\textit{\color{bookprimary}``The benefits of the phone are real. But many people use a few grains of sugar to excuse the pile of garbage they swallow every day.''}
\end{center}

\ngancach

\begin{tcolorbox}[colback=booksecondary!8,colframe=booksecondary!40,arc=5pt,title={\bfseries Góc Suy Nghĩ},fonttitle=\bfseries\color{booksecondary}]
Điện thoại giúp được mười việc tốt.
\textit{(A phone can help with ten good things.)}

Nhưng rất nhiều người dùng nó cho một trăm việc thừa.
\textit{(But many people use it for a hundred unnecessary things.)}

Tra cứu tài liệu là ưu điểm.
\textit{(Searching for learning material is an advantage.)}

Lướt vô tận tới nát đầu thì không.
\textit{(Scrolling endlessly until the mind is wrecked is not.)}

Gọi người thân khi khẩn cấp là ưu điểm.
\textit{(Calling loved ones in an emergency is an advantage.)}

Để một thông báo vô nghĩa giật mình mình suốt ngày thì không.
\textit{(Letting pointless notifications yank your nerves all day is not.)}

Vấn đề không phải điện thoại không có ích.
\textit{(The problem is not that the phone lacks usefulness.)}

Vấn đề là cái ít ích đang bị dùng như giấy phép cho quá nhiều điều hại.
\textit{(The problem is that small usefulness is being used as a permit for too much harm.)}
\end{tcolorbox}

\section{Một Chút Lợi Không Rửa Sạch Một Đống Hại}
\begin{truyen}{Một Chút Lợi Không Rửa Sạch Một Đống Hại}{Small Benefits Do Not Wash Away Big Damage}
\chuhoa{C}{ó bạn nói rất chắc}: ``Nhưng điện thoại giúp em học tiếng Anh, xem bài giảng, tra từ điển mà thầy.''
\textit{(A student says very confidently: ``But the phone helps me learn English, watch lessons, and look up vocabulary, teacher.'')}

Thầy hỏi luôn: ``Mỗi ngày em dùng được bao lâu cho mấy việc đó?''
\textit{(The teacher asks at once: ``How much of your day actually goes to those things?'')}

Bạn ấy ấp úng.
\textit{(The student hesitates.)}

\teacherpair{Nếu em dùng mười lăm phút để học mà dùng ba tiếng để lướt vô định, thì đừng lấy mười lăm phút sạch sẽ đó ra xin giảm án cho ba tiếng bầy hầy kia. Con người rất hay tự lừa kiểu này: bấu vào một chút điều tử tế mình có làm để khỏi phải nhìn thẳng vào phần mình đang tự phá mình nhiều hơn.}{``If you spend fifteen minutes learning but three hours wandering aimlessly, do not pull out those fifteen clean minutes to ask pardon for the three messy hours. People deceive themselves this way all the time: they cling to one small decent thing they did so they can avoid facing the much larger part where they are damaging themselves.''}

Thầy nói tiếp: ``Ưu điểm phải được nhìn bằng tỷ lệ. Một cái có ích ít mà bị bọc trong quá nhiều cái hại thì cuối cùng tổng đời sống vẫn đi xuống. Đừng lấy một giọt thuốc nhỏ vào ly nước bẩn rồi bảo thế là ly đó đáng uống.''
\textit{(The teacher continues: ``Strengths must be judged by proportion. If a small benefit is wrapped inside too much harm, then the total direction of life still goes downward. Do not drip one drop of medicine into a dirty glass of water and then declare it safe to drink.'')}
\end{truyen}

\section{Năm Phút Học, Năm Mươi Phút Lướt}
\begin{truyen}{Năm Phút Học, Năm Mươi Phút Lướt}{Five Minutes of Learning, Fifty Minutes of Drift}
\chuhoa{C}{ó người mở điện thoại để tra đúng một khái niệm.}
\textit{(Someone opens the phone to search for one concept.)}

Một tiếng sau mới giật mình vì đã đi qua đủ video, bình luận và tin rác.
\textit{(An hour later they suddenly realize they have drifted through videos, comments, and junk news.)}

\teacherpair{Đó là kiểu mất đời rất thời thượng: khởi đầu bằng một lý do tử tế rồi kết thúc bằng một đống phân tán. Các em rất thích kể cái mở đầu đẹp. Nhưng đời các em bị tính bằng chỗ kết thúc thật.}{``That is a fashionable way of losing life: beginning with a decent reason and ending in a pile of fragmentation. Students love telling the beautiful opening. But life is measured by the actual ending point.''}
\end{truyen}

\section{Một Bài Giảng Không Cứu Nổi Một Đêm Rác}
\begin{truyen}{Một Bài Giảng Không Cứu Nổi Một Đêm Rác}{One Lecture Does Not Redeem a Night of Trash}
\chuhoa{C}{ó bạn xem được một video bài giảng rất hay rồi tự thấy mình chăm.}
\textit{(A student watches one excellent lecture video and feels diligent.)}

Nhưng sau đó lại đổ cả đêm vào thứ nội dung làm đầu óc mục thêm.
\textit{(But then spends the whole night on content that rots the mind further.)}

\teacherpair{Con người rất thích tự thưởng huân chương sớm. Làm được một điều đúng không tự động xóa chín điều đang làm sai. Nếu em ăn một muỗng rau rồi nuốt tiếp cả mâm rác, bụng em không nhờ muỗng rau đó mà sạch hơn.}{``People love giving themselves medals too early. Doing one thing right does not erase nine things being done wrong. If you eat one spoon of vegetables and then swallow a tray of junk, your body will not be purified by that spoon of greens.''}
\end{truyen}

\section{Ngụy Biện Bằng Tiếng Anh}
\begin{truyen}{Ngụy Biện Bằng Tiếng Anh}{Excusing Yourself with English Learning}
\chuhoa{N}{hiều bạn nói}: ``Em lướt TikTok chủ yếu để học tiếng Anh.''
\textit{(Many students say: ``I mainly scroll TikTok to learn English.'')}

Thầy hỏi: ``Ngoài mấy câu ngắn em nhặt được, kỹ năng nào của em thật sự tăng thấy rõ?''
\textit{(The teacher asks: ``Besides the short phrases you picked up, which skill of yours has actually improved in a visible way?'')}

Bạn ấy đứng lại.
\textit{(The student pauses.)}

\teacherpair{Học tiếng Anh là một ưu điểm thật. Nhưng mượn tiếng Anh để rửa mặt cho sự nghiện lướt thì rất bẩn. Thứ gì thật sự giúp em mạnh lên phải để lại dấu vết rõ trong vốn từ, cách nghe, cách nói, cách viết, chứ không chỉ để lại vài câu nói cho vui mồm.}{``Learning English is a real advantage. But borrowing English to wash the face of a scrolling addiction is dirty work. Whatever truly makes you stronger must leave clear traces in your vocabulary, listening, speaking, and writing, not merely in a few catchy lines for amusement.''}
\end{truyen}

\section{Nói Làm Việc Nhưng Toàn Chuyển Cảnh}
\begin{truyen}{Nói Làm Việc Nhưng Toàn Chuyển Cảnh}{Calling It Work While Constantly Switching Scenes}
\chuhoa{C}{ó người mở điện thoại vì công việc.}
\textit{(Some people open the phone for work.)}

Nhưng cứ vài phút lại trôi sang một ứng dụng khác.
\textit{(But every few minutes they slide into another app.)}

\teacherpair{Làm việc bằng điện thoại không sai. Sai là cứ lấy chữ công việc làm cái áo đẹp cho một đầu óc phân tán. Nhiều người tự nhận bận công việc, thật ra chỉ đang bận bị màn hình giật qua giật lại.}{``Working through a phone is not wrong. What is wrong is using the word work as a beautiful coat for a scattered mind. Many people claim to be busy with work when in truth they are merely busy being jerked around by a screen.''}
\end{truyen}

\section{Một Cuộc Gọi Tử Tế Không Cứu Nổi Cả Ngày Cáu Bẳn}
\begin{truyen}{Một Cuộc Gọi Tử Tế Không Cứu Nổi Cả Ngày Cáu Bẳn}{One Useful Call Does Not Redeem a Whole Irritated Day}
\chuhoa{C}{ó người nói điện thoại rất cần vì nhờ nó họ gọi được việc cần.}
\textit{(Someone says the phone is necessary because it helps them make useful calls.)}

Đúng.
\textit{(True.)}

Nhưng nếu phần còn lại của ngày là bị thông báo giật cho nhấp nhổm, cáu gắt và mệt thần kinh thì cái gọi là tiện kia cũng đang bị mua bằng giá quá đắt.
\textit{(But if the rest of the day is spent being jerked around by notifications into restlessness, irritability, and nervous fatigue, then that convenience is being bought at far too high a price.)}

\teacherpair{Đừng chỉ hỏi điện thoại cho mình cái gì. Hãy hỏi nó đang lấy của mình cái gì. Người lớn lên là người biết tính giá, không phải chỉ biết kể ưu điểm.}{``Do not only ask what the phone gives you. Ask what it is taking from you. Growing up means learning to calculate costs, not merely listing benefits.''}
\end{truyen}

\section{Cái Tiện Làm Hỏng Nhiều Thói Quen Khó}
\begin{truyen}{Cái Tiện Làm Hỏng Nhiều Thói Quen Khó}{Convenience Weakens Hard Habits}
\chuhoa{M}{ột thời, con người phải chịu khó hơn để tìm câu trả lời, để hẹn nhau, để chờ nhau.}
\textit{(There was a time when people had to work harder to find answers, make plans, and wait for one another.)}

Bây giờ cái tiện đến quá nhanh.
\textit{(Now convenience arrives too quickly.)}

\teacherpair{Cái tiện là một phước lớn nếu người cầm nó có kỷ luật. Nhưng cái tiện rơi vào tay một người lười tự chủ thì rất dễ làm hỏng những cơ bắp tinh thần như nhẫn nại, tập trung, chờ đợi và làm đến nơi đến chốn.}{``Convenience is a great blessing if the person holding it has discipline. But when convenience falls into the hands of someone lacking self-control, it easily weakens the mental muscles of patience, focus, waiting, and finishing things properly.''}
\end{truyen}

\section{Càng Hữu Ích Càng Phải Giữ Giới Hạn}
\begin{truyen}{Càng Hữu Ích Càng Phải Giữ Giới Hạn}{The More Useful It Is, the More Limits It Needs}
\chuhoa{N}{gười ta thường đặt giới hạn với thứ mình biết là xấu rõ ràng.}
\textit{(People usually set limits on things they clearly know are bad.)}

Còn với thứ hữu ích, họ dễ mở cửa quá rộng.
\textit{(But with something useful, they easily open the door too wide.)}

\teacherpair{Đó là cái bẫy lớn của điện thoại. Nó không xấu thô như thuốc độc. Nó tốt vừa đủ để được tha thứ mỗi ngày. Thành ra muốn sống tỉnh, em càng phải biết đặt giới hạn với thứ có ích. Vì chính thứ có ích mới dễ lấn sâu nhất.}{``That is the major trap of the phone. It is not crudely evil like poison. It is useful enough to earn daily forgiveness. That is why, if you want to live awake, you must set limits precisely around useful things. Useful things are often the ones that invade deepest.''}
\end{truyen}

\section{Tỷ Lệ Mới Là Bản Án}
\begin{truyen}{Tỷ Lệ Mới Là Bản Án}{The Ratio Is the Verdict}
\chuhoa{M}{ột điều có ích không vô nghĩa.}
\textit{(A useful action is not meaningless.)}

Nhưng giá trị thật của nó nằm trong tổng thể.
\textit{(But its real value lies in the total pattern.)}

\teacherpair{Đừng hỏi điện thoại của em có giúp được gì không. Hỏi tỷ lệ. Tỷ lệ dùng để học, để làm việc, để kết nối tử tế so với tỷ lệ lướt rác, né đời, cắt mạch, mất ngủ và mỏi đầu là bao nhiêu. Tỷ lệ đó mới là bản án công bằng nhất cho thói quen của em.}{``Do not ask whether your phone helps with anything. Ask for the ratio. What is the proportion of learning, work, and decent connection compared to junk scrolling, life avoidance, broken focus, lost sleep, and mental fatigue? That ratio is the fairest verdict on your habits.''}
\end{truyen}

\section{Tỉnh Ngộ Không Bắt Đầu Từ Việc Vứt Máy}
\begin{truyen}{Tỉnh Ngộ Không Bắt Đầu Từ Việc Vứt Máy}{Awakening Does Not Begin by Throwing the Phone Away}
\chuhoa{N}{hiều người nghe phê phán điện thoại xong lại nghĩ giải pháp là cực đoan: bỏ sạch, cắt sạch, ghét sạch.}
\textit{(Many people hear criticism of phones and immediately imagine an extreme solution: abandon everything, cut everything, hate everything.)}

Không.
\textit{(No.)}

\teacherpair{Tỉnh ngộ bắt đầu từ chỗ tính sòng phẳng. Điện thoại có cái lợi nào, giữ. Cái hại nào đang nuốt mình, chặt. Người trưởng thành không cần đóng kịch thanh cao. Họ chỉ cần đủ tỉnh để không tiếp tục tự bào chữa.}{``Awakening begins with honest accounting. Whatever benefit the phone truly brings, keep it. Whatever harm is swallowing you, cut it off. A mature person does not need to perform false nobility. They only need enough clarity to stop excusing themselves.''}
\end{truyen}

\begin{baihoc}
Điện thoại có ưu điểm.
\textit{(The phone has advantages.)}

Nhưng ở rất nhiều người, ưu chỉ là chút lý do đẹp để che một đời sống đang bị cắt vụn, mỏi mệt và nông dần đi.
\textit{(But in many people, those strengths are only pretty excuses covering a life that is becoming fragmented, tired, and increasingly shallow.)}
\end{baihoc}

\begin{tuhocnhanh}
1. Trong một ngày thật của em, tỉ lệ dùng điện thoại cho việc có ích và vô ích là bao nhiêu? \textit{(In your real day, what is the ratio between useful and useless phone use?)}

2. Em có đang lấy một ít điều tốt để che một đống điều tệ không? \textit{(Are you using a few good things to cover a pile of bad ones?)}

3. Nếu nhìn công bằng, điện thoại đang kéo tổng đời em đi lên hay đi xuống? \textit{(If judged fairly, is the phone pulling your overall life upward or downward?)}
\end{tuhocnhanh}

\begin{tongketchuong}
Vấn đề của nhiều người không phải là không biết điện thoại có lợi.
\textit{(Many people do not suffer from ignorance of the phone's benefits.)}

Vấn đề là họ cứ lôi cái lợi ít ỏi đó ra để che cho một đời sống bị tiêu hao quá nhiều.
\textit{(The problem is that they keep dragging out those small benefits to hide a life being consumed far too much.)}
\end{tongketchuong}

\ngancach